\section*{Appendices}
\phantomsection
\addcontentsline{toc}{section}{Appendices}

\subsection*{Appendix A. Detailed Data-Quality Results}
\phantomsection
\addcontentsline{toc}{subsection}{Appendix A. Detailed Data-Quality Results}

The eight annual files hold 48,538,713 observations from 91 station IDs. Every record is a five-minute rainfall total in mm. No missing timestamp, station, rainfall, coordinate or update-time value, no negative or non-finite rainfall, and no duplicate original (station, timestamp) key was found in any year.

\begingroup
\small
\setlength{\parskip}{0pt}
\setlength{\tabcolsep}{3pt}
\renewcommand{\arraystretch}{1.16}
\begin{longtable}{@{}>{\raggedright\arraybackslash}p{(\linewidth-8\tabcolsep)*\real{0.09}}>{\raggedright\arraybackslash}p{(\linewidth-8\tabcolsep)*\real{0.14}}>{\raggedright\arraybackslash}p{(\linewidth-8\tabcolsep)*\real{0.12}}>{\raggedright\arraybackslash}p{(\linewidth-8\tabcolsep)*\real{0.325}}>{\raggedright\arraybackslash}p{(\linewidth-8\tabcolsep)*\real{0.325}}@{}}
\caption*{\textbf{Table A.1.} Dataset overview by year}\\
\toprule
\textbf{Year} & \textbf{Rows} & \textbf{Stations} & \textbf{First observation} & \textbf{Last observation} \\
\midrule
\endfirsthead
\toprule
\textbf{Year} & \textbf{Rows} & \textbf{Stations} & \textbf{First observation} & \textbf{Last observation} \\
\midrule
\endhead
2017 & 5,256,106 & 62 & 2017-01-01 08:04:59 & 2017-12-31 23:55:00 \\
2018 & 4,594,246 & 56 & 2018-01-01 00:00:00 & 2018-12-31 23:55:00 \\
2019 & 5,026,862 & 54 & 2019-01-01 00:00:00 & 2019-12-31 23:55:00 \\
2020 & 6,334,473 & 80 & 2020-01-01 00:00:00 & 2020-12-31 23:55:00 \\
2021 & 7,044,646 & 76 & 2021-01-01 00:00:00 & 2021-12-31 23:55:00 \\
2022 & 7,133,230 & 72 & 2022-01-01 00:00:00 & 2022-12-31 23:55:00 \\
2023 & 6,742,890 & 74 & 2023-01-01 00:00:00 & 2023-12-31 23:55:00 \\
2024 & 6,406,260 & 70 & 2024-01-01 00:00:00 & 2024-12-31 23:55:00 \\
\textbf{Total} & \textbf{48,538,713} & \textbf{91 unique} &  &  \\
\bottomrule
\end{longtable}
\endgroup

Coverage is the share of the expected five-minute records that a station reported between its own first and last observation in that year, so a station that joined or left during the year is not penalised. It differs from the calendar-based coverage of Section 3.3 and Figure 1, which counts against the full month or year. The audit coverage is diagnostic: it was not used to delete any station.

\begingroup
\small
\setlength{\parskip}{0pt}
\setlength{\tabcolsep}{3pt}
\renewcommand{\arraystretch}{1.16}
\begin{longtable}{@{}>{\raggedright\arraybackslash}p{(\linewidth-10\tabcolsep)*\real{0.08}}>{\raggedright\arraybackslash}p{(\linewidth-10\tabcolsep)*\real{0.12}}>{\raggedright\arraybackslash}p{(\linewidth-10\tabcolsep)*\real{0.17}}>{\raggedright\arraybackslash}p{(\linewidth-10\tabcolsep)*\real{0.17}}>{\raggedright\arraybackslash}p{(\linewidth-10\tabcolsep)*\real{0.23}}>{\raggedright\arraybackslash}p{(\linewidth-10\tabcolsep)*\real{0.23}}@{}}
\caption*{\textbf{Table A.2.} Station coverage by year}\\
\toprule
\textbf{Year} & \textbf{Stations} & \textbf{Mean coverage} & \textbf{Median coverage} & \textbf{Stations below 80\%} & \textbf{Stations below 50\%} \\
\midrule
\endfirsthead
\toprule
\textbf{Year} & \textbf{Stations} & \textbf{Mean coverage} & \textbf{Median coverage} & \textbf{Stations below 80\%} & \textbf{Stations below 50\%} \\
\midrule
\endhead
2017 & 62 & 0.8972 & 0.9497 & 5 & 3 \\
2018 & 56 & 0.8110 & 0.8179 & 17 & 0 \\
2019 & 54 & 0.9030 & 0.9490 & 8 & 2 \\
2020 & 80 & 0.9144 & 0.9481 & 5 & 2 \\
2021 & 76 & 0.9563 & 0.9842 & 3 & 1 \\
2022 & 72 & 0.9775 & 0.9975 & 2 & 1 \\
2023 & 74 & 0.9331 & 0.9677 & 2 & 2 \\
2024 & 70 & 0.9608 & 0.9835 & 2 & 0 \\
\bottomrule
\end{longtable}
\endgroup

Source: \path{task1/results/data_quality_audit/yearly_quality.csv} and \path{coverage_by_year.csv}, written by \path{task1/code/rainfall_audit.py}; discussion in \path{task1/docs/data_quality_findings.md}.

\clearpage
\subsection*{Appendix B. Full Feature Definitions}
\phantomsection
\addcontentsline{toc}{subsection}{Appendix B. Full Feature Definitions}

\begingroup
\small
\setlength{\parskip}{0pt}
\setlength{\tabcolsep}{3pt}
\renewcommand{\arraystretch}{1.16}
\begin{longtable}{@{}>{\raggedright\arraybackslash}p{(\linewidth-8\tabcolsep)*\real{0.055}}>{\raggedright\arraybackslash}p{(\linewidth-8\tabcolsep)*\real{0.255}}>{\raggedright\arraybackslash}p{(\linewidth-8\tabcolsep)*\real{0.05}}>{\raggedright\arraybackslash}p{(\linewidth-8\tabcolsep)*\real{0.43}}>{\raggedright\arraybackslash}p{(\linewidth-8\tabcolsep)*\real{0.21}}@{}}
\caption*{\textbf{Table B.1.} Features}\\
\toprule
\textbf{Tier} & \textbf{Feature(s)} & \textbf{No.} & \textbf{How it is constructed} & \textbf{Data used} \\
\midrule
\endfirsthead
\toprule
\textbf{Tier} & \textbf{Feature(s)} & \textbf{No.} & \textbf{How it is constructed} & \textbf{Data used} \\
\midrule
\endhead
T0 & \path{t0_x_km}, \path{t0_y_km} & 2 & Station position in km east and north of the island centre (103.82\(^{\circ}\)E, 1.35\(^{\circ}\)N), by an equirectangular projection & Latest station coordinates \\
T0 & \path{t0_wet_ratio} & 1 & (O + 500) / (E + 500): rainy 30-minute windows observed at the station (Y \(>\) 0), over the number expected if it followed the network at the same times; shrunk towards 1 & Labels 2017--2022, constant per station \\
T0 & \path{t0_heavy_ratio} & 1 & The same ratio for heavy windows (Y \(>\) 5 mm) & Labels 2017--2022, constant per station \\
T0 & \path{t0_hod_s1}--\path{s3}, \path{t0_hod_c1}--\path{c3} & 6 & sin and cos of 2\(\pi\)k\(\cdot\)h/24 for k = 1, 2, 3, where h is the local hour of the prediction time (with minutes) & Prediction time \\
T0 & \path{t0_doy_s1}, \path{s2}, \path{t0_doy_c1}, \path{c2} & 4 & sin and cos of 2\(\pi\)k\(\cdot\)\(\delta\) for k = 1, 2, where \(\delta\) is the fraction of the year elapsed & Prediction time \\
T0 & \path{t0_season} & 1 & Monsoon code: 0 for Dec--Mar, 2 for Jun--Sep, 1 otherwise & Prediction time \\
T1 & \path{t1_rain_5}, \path{t1_rain_15}, \path{t1_rain_30}, \path{t1_rain_60}, \path{t1_rain_120}, \path{t1_rain_1440} & 6 & Rain total at the station over the last 5, 15, 30, 60, 120 and 1,440 minutes (\path{t1_rain_5} is the current reading) & Station's own readings, t and earlier \\
T1 & \path{t1_max5_30} & 1 & Largest five-minute reading in the last 30 minutes & Same \\
T1 & \path{t1_wet_slots_60} & 1 & Number of wet readings among the last 12 & Same \\
T1 & \path{t1_trend_30} & 1 & 2\(\cdot\)\(R^{15}\) \(-\) \(R^{30}\): the last 15 minutes minus the 15 minutes before, positive when rain is intensifying & Same \\
T1 & \path{t1_min_since_rain} & 1 & Minutes since the last wet reading, capped at 1,440 & Same \\
T1 & \path{t1_missing_120} & 1 & Number of missing readings among the last 24 & Same \\
T2 & \path{t2_idw_rain_15}, \path{t2_idw_rain_30} & 2 & Mean 15- and 30-minute rain of the other stations within 10 km, weighted by \(1/\max(d,0.5)^2\), reporting stations only (0 if none) & Other stations, time t \\
T2 & \path{t2_idw_trend} & 1 & 2 \(\times\) \path{t2_idw_rain_15} \(-\) \path{t2_idw_rain_30}, the neighbours' version of the trend & Same \\
T2 & \path{t2_far_rain_30} & 1 & The same weighted mean of 30-minute rain for stations 10 to 25 km away & Same \\
T2 & \path{t2_max_rain_30} & 1 & Largest 30-minute rain at any reporting station within 10 km & Same \\
T2 & \path{t2_wet_frac_5km}, \path{t2_wet_frac_10km} & 2 & Share of reporting stations within 5 km and 10 km with rain in the last 30 minutes & Same \\
T2 & \path{t2_dist_rain_km} & 1 & Distance to the nearest station with rain in the last 15 minutes, capped at 30 km & Same \\
T2 & \path{t2_rain_30_n}, \path{_e}, \path{_s}, \path{_w} & 4 & Mean 30-minute rain of the reporting stations within 25 km in the north, east, south and west quadrants & Same \\
T2 & \path{t2_island_wet_frac}, \path{t2_island_rain_30} & 2 & Share of all reporting stations that are wet, and their mean 30-minute rain & All stations, time t \\
T2 & \path{t2_nbrs_10km} & 1 & Number of reporting stations within 10 km (tells the model how many stations the averages rest on) & Other stations, time t \\
T3 & \path{t3_d15_idw_rain_30}, \path{_far_rain_30}, \path{_max_rain_30}, \path{_wet_frac_5km}, \path{_wet_frac_10km}, \path{_dist_rain_km}, \path{_rain_30_n}, \path{_e}, \path{_s}, \path{_w}, \path{_island_wet_frac}, \path{_island_rain_30} & 12 & Each of these twelve Tier 2 summaries now, minus the same summary recomputed from the network readings 15 minutes earlier & Network at t and at t \(-\) 15 min \\
T3 & \path{t3_up15_rain}, \path{t3_up15_wet} & 2 & Gaussian-weighted mean (\(\sigma\) = 4 km, cut-off 10 km) of 15-minute rain, and of the wet indicator, around the point 15 minutes upstream: the station moved backwards along the seasonal storm motion (Table 5.3) & Network at t \\
T3 & \path{t3_up30_rain}, \path{t3_up30_wet} & 2 & The same around the point 30 minutes upstream & Network at t \\
T3 & \path{t3_up30_nbrs} & 1 & Number of stations near the 30-minute upstream point (it separates "dry" from "no gauge there") & Network at t \\
\bottomrule
\end{longtable}
\endgroup

\clearpage
\subsection*{Appendix C. Hyperparameter Tuning and Test Performance}
\phantomsection
\addcontentsline{toc}{subsection}{Appendix C. Hyperparameter Tuning and Test Performance}

\begingroup
\small
\setlength{\parskip}{0pt}
\setlength{\tabcolsep}{3pt}
\renewcommand{\arraystretch}{1.16}
\begin{longtable}{@{}>{\raggedright\arraybackslash}p{(\linewidth-10\tabcolsep)*\real{0.2}}>{\raggedright\arraybackslash}p{(\linewidth-10\tabcolsep)*\real{0.15}}>{\raggedright\arraybackslash}p{(\linewidth-10\tabcolsep)*\real{0.18}}>{\raggedright\arraybackslash}p{(\linewidth-10\tabcolsep)*\real{0.17}}>{\raggedright\arraybackslash}p{(\linewidth-10\tabcolsep)*\real{0.13}}>{\raggedright\arraybackslash}p{(\linewidth-10\tabcolsep)*\real{0.17}}@{}}
\caption*{\textbf{Table C.1.} Logistic regression on the 2023 validation set}\\
\toprule
\textbf{Tiers} & \textbf{\(\lambda\) (regParam)} & \textbf{Non-linear inputs} & \textbf{Validation RPS} & \textbf{Iterations} & \textbf{Fit (min)} \\
\midrule
\endfirsthead
\toprule
\textbf{Tiers} & \textbf{\(\lambda\) (regParam)} & \textbf{Non-linear inputs} & \textbf{Validation RPS} & \textbf{Iterations} & \textbf{Fit (min)} \\
\midrule
\endhead
T0+T1+T2 & 1e-05 & no & 0.0961 & 429 & 4.7 \\
T0+T1+T2 & 1e-06 & yes & 0.0935 & 659 & 10.7 \\
T0+T1+T2 & 1e-05 & yes & 0.0936 & 609 & 9.9 \\
T0+T1+T2 & 0.0001 & yes & 0.0937 & 530 & 8.8 \\
T0+T1+T2+T3 & 1e-06 & yes & 0.0922 & 655 & 12.9 \\
\bottomrule
\end{longtable}
\endgroup

The first row is the plain model; the last row adds Tier 3 and is the final specification.

Rows are in the order tried; the last row is the chosen setting before Tier 3 was added.

\begingroup
\footnotesize
\setlength{\parskip}{0pt}
\setlength{\tabcolsep}{3pt}
\renewcommand{\arraystretch}{1.16}
\begin{longtable}{@{}>{\raggedright\arraybackslash}p{(\linewidth-14\tabcolsep)*\real{0.1}}>{\raggedright\arraybackslash}p{(\linewidth-14\tabcolsep)*\real{0.14}}>{\raggedright\arraybackslash}p{(\linewidth-14\tabcolsep)*\real{0.2}}>{\raggedright\arraybackslash}p{(\linewidth-14\tabcolsep)*\real{0.1}}>{\raggedright\arraybackslash}p{(\linewidth-14\tabcolsep)*\real{0.1}}>{\raggedright\arraybackslash}p{(\linewidth-14\tabcolsep)*\real{0.13}}>{\raggedright\arraybackslash}p{(\linewidth-14\tabcolsep)*\real{0.11}}>{\raggedright\arraybackslash}p{(\linewidth-14\tabcolsep)*\real{0.12}}@{}}
\caption*{\textbf{Table C.2.} Random forest screening on 2023}\\
\toprule
\textbf{max\newline Depth} & \textbf{min\newline Instances\newline PerNode} & \textbf{featureSubset\newline Strategy} & \textbf{maxBins} & \textbf{Impurity} & \textbf{Validation RPS} & \textbf{Fit (min)} & \textbf{Nodes} \\
\midrule
\endfirsthead
\toprule
\textbf{max\newline Depth} & \textbf{min\newline Instances\newline PerNode} & \textbf{featureSubset\newline Strategy} & \textbf{maxBins} & \textbf{Impurity} & \textbf{Validation RPS} & \textbf{Fit (min)} & \textbf{Nodes} \\
\midrule
\endhead
14 & 100 & sqrt & 64 & gini & 0.0927 & 2.7 & 224,268 \\
17 & 100 & sqrt & 64 & gini & 0.0923 & 4.5 & 379,530 \\
20 & 100 & sqrt & 64 & gini & 0.0923 & 6.0 & 435,878 \\
24 & 100 & sqrt & 64 & gini & 0.0923 & 7.3 & 444,966 \\
24 & 20 & sqrt & 64 & gini & 0.0930 & 21.5 & 2,675,066 \\
24 & 50 & sqrt & 64 & gini & 0.0920 & 11.4 & 984,782 \\
24 & 300 & sqrt & 64 & gini & 0.0935 & 4.0 & 122,922 \\
24 & 50 & onethird & 64 & gini & 0.0919 & 20.3 & 1,014,978 \\
24 & 50 & 0.5 & 64 & gini & 0.0919 & 27.8 & 1,029,336 \\
24 & 50 & onethird & 256 & gini & 0.0912 & 56.1 & 1,030,260 \\
24 & 50 & onethird & 256 & entropy & 0.0911 & 63.8 & 1,044,176 \\
\bottomrule
\end{longtable}
\endgroup

RPS: ranked probability score; Brier \(>\)\(\theta\): Brier score of P(Y \(>\) \(\theta\)) in mm; AUC-PR: area under the precision-recall curve. Lower is better for RPS, log loss and Brier; higher for AUC-PR.

\begin{landscape}
\begingroup
\footnotesize
\setlength{\parskip}{0pt}
\setlength{\tabcolsep}{3pt}
\renewcommand{\arraystretch}{1.16}
\begin{longtable}{@{}>{\raggedright\arraybackslash}p{(\linewidth-22\tabcolsep)*\real{0.115}}>{\raggedright\arraybackslash}p{(\linewidth-22\tabcolsep)*\real{0.15}}>{\raggedright\arraybackslash}p{(\linewidth-22\tabcolsep)*\real{0.0735}}>{\raggedright\arraybackslash}p{(\linewidth-22\tabcolsep)*\real{0.0735}}>{\raggedright\arraybackslash}p{(\linewidth-22\tabcolsep)*\real{0.0735}}>{\raggedright\arraybackslash}p{(\linewidth-22\tabcolsep)*\real{0.0735}}>{\raggedright\arraybackslash}p{(\linewidth-22\tabcolsep)*\real{0.0735}}>{\raggedright\arraybackslash}p{(\linewidth-22\tabcolsep)*\real{0.0735}}>{\raggedright\arraybackslash}p{(\linewidth-22\tabcolsep)*\real{0.0735}}>{\raggedright\arraybackslash}p{(\linewidth-22\tabcolsep)*\real{0.0735}}>{\raggedright\arraybackslash}p{(\linewidth-22\tabcolsep)*\real{0.0735}}>{\raggedright\arraybackslash}p{(\linewidth-22\tabcolsep)*\real{0.0735}}@{}}
\caption*{\textbf{Table C.3.} Test year 2024}\\
\toprule
\textbf{Subset} & \textbf{Forecast} & \textbf{RPS} & \textbf{log loss} & \textbf{Brier \(>\)0} & \textbf{Brier \(>\)1} & \textbf{Brier \(>\)5} & \textbf{Brier \(>\)20} & \textbf{AUC-PR \(>\)0} & \textbf{AUC-PR \(>\)1} & \textbf{AUC-PR \(>\)5} & \textbf{AUC-PR \(>\)20} \\
\midrule
\endfirsthead
\toprule
\textbf{Subset} & \textbf{Forecast} & \textbf{RPS} & \textbf{log loss} & \textbf{Brier \(>\)0} & \textbf{Brier \(>\)1} & \textbf{Brier \(>\)5} & \textbf{Brier \(>\)20} & \textbf{AUC-PR \(>\)0} & \textbf{AUC-PR \(>\)1} & \textbf{AUC-PR \(>\)5} & \textbf{AUC-PR \(>\)20} \\
\midrule
\endhead
All forecasts & Constant & 0.1903 & 0.3545 & 0.0596 & 0.0244 & 0.00825 & 0.00103 & 0.0636 & 0.0250 & 0.0083 & 0.0010 \\
All forecasts & Climatology & 0.1864 & 0.3437 & 0.0578 & 0.0240 & 0.00820 & 0.00102 & 0.1254 & 0.0540 & 0.0234 & 0.0043 \\
All forecasts & Persistence & 0.1246 & 0.2577 & 0.0355 & 0.0170 & 0.00701 & 0.00100 & 0.6406 & 0.4731 & 0.2500 & 0.0495 \\
All forecasts & Logistic regression & 0.0922 & 0.1933 & 0.0259 & 0.0126 & 0.00547 & 0.00087 & 0.7850 & 0.6888 & 0.5066 & 0.1956 \\
All forecasts & Random forest & 0.0887 & 0.1863 & 0.0248 & 0.0121 & 0.00539 & 0.00088 & 0.8015 & 0.7117 & 0.5246 & 0.2285 \\
Dry now & Constant & 0.0708 & 0.1760 & 0.0260 & 0.0081 & 0.00275 & 0.00026 & 0.0253 & 0.0079 & 0.0027 & 0.0003 \\
Dry now & Climatology & 0.0710 & 0.1705 & 0.0259 & 0.0082 & 0.00277 & 0.00026 & 0.0522 & 0.0208 & 0.0084 & 0.0009 \\
Dry now & Persistence & 0.0679 & 0.1611 & 0.0247 & 0.0079 & 0.00272 & 0.00026 & 0.0253 & 0.0079 & 0.0027 & 0.0003 \\
Dry now & Logistic regression & 0.0544 & 0.1142 & 0.0187 & 0.0066 & 0.00243 & 0.00025 & 0.3962 & 0.2608 & 0.1616 & 0.0410 \\
Dry now & Random forest & 0.0530 & 0.1111 & 0.0182 & 0.0065 & 0.00241 & 0.00025 & 0.4298 & 0.2959 & 0.1903 & 0.0553 \\
Raining now & Constant & 1.9499 & 2.9833 & 0.5544 & 0.2635 & 0.08926 & 0.01232 & 0.6275 & 0.2759 & 0.0907 & 0.0123 \\
Raining now & Climatology & 1.8859 & 2.8948 & 0.5278 & 0.2578 & 0.08825 & 0.01228 & 0.6403 & 0.2977 & 0.1204 & 0.0236 \\
Raining now & Persistence & 0.9580 & 1.6771 & 0.1935 & 0.1515 & 0.07024 & 0.01185 & 0.8230 & 0.6128 & 0.3401 & 0.0628 \\
Raining now & Logistic regression & 0.6487 & 1.3581 & 0.1324 & 0.1004 & 0.05017 & 0.01001 & 0.9324 & 0.8222 & 0.6352 & 0.2974 \\
Raining now & Random forest & 0.6140 & 1.2924 & 0.1219 & 0.0952 & 0.04934 & 0.01009 & 0.9427 & 0.8418 & 0.6466 & 0.2885 \\
\bottomrule
\end{longtable}
\endgroup
\end{landscape}

\clearpage
\subsection*{Appendix D. Reproduction Instructions and Model Files}
\phantomsection
\addcontentsline{toc}{subsection}{Appendix D. Reproduction Instructions and Model Files}

The code, documentation, trained models and small result tables are in the repository (\url{https://github.com/LuweiChenbot/RainfallPatternPrediction}). Raw data and the generated Parquet tables are not included.

\subsubsection*{D.1 Environment}
\phantomsection
\addcontentsline{toc}{subsubsection}{D.1 Environment}

Java 17 or 21, Python 3.12, PySpark 4.2.0. From the repository root:

\begin{verbatim}
python3 -m venv .venv
source .venv/bin/activate
pip install -r requirements.txt
cp .env.example .env     # set cores and memory for the machine
\end{verbatim}

\subsubsection*{D.2 Run Steps}
\phantomsection
\addcontentsline{toc}{subsubsection}{D.2 Run Steps}

\begin{enumerate}
\item Download the eight annual files of the \textit{Historical Rainfall across Singapore} collection (data.gov.sg) into \path{data/}.
\item Task 1: run, in order, \path{task1/code/rainfall_audit.py} (optional), \path{prepare_base_data.py}, \path{prepare_temporal_data.py}, \path{fetch_boundary.py} and \path{build_matrix.py}, then the four \path{plot_*.py} scripts (commands in \path{task1/README.md}). This produces the five-minute station matrix \path{task1/data/matrix_5min/} and the Task 1 figures.
\item Task 2: \texttt{python task2/run\_pipeline.py} runs notebooks 01, 02, 05, 06, 08, 03, 04, 07, 09 and 10 in that order. Settings of the final run: \path{RAIN_TRAIN=train_full}, \path{RAIN_RF_TREES=50}, \path{RAIN_RF_MEMORY_MB=512}.
\item Figures and the replay page: \texttt{python task2/code/plot\_results.py} and \texttt{python task2/code/build\_replay.py}.
\end{enumerate}

\subsubsection*{D.3 Submitted model files}
\phantomsection
\addcontentsline{toc}{subsubsection}{D.3 Submitted model files}

\begingroup
\small
\setlength{\parskip}{0pt}
\setlength{\tabcolsep}{3pt}
\renewcommand{\arraystretch}{1.16}
\begin{longtable}{@{}>{\raggedright\arraybackslash}p{(\linewidth-4\tabcolsep)*\real{0.35}}>{\raggedright\arraybackslash}p{(\linewidth-4\tabcolsep)*\real{0.53}}>{\raggedright\arraybackslash}p{(\linewidth-4\tabcolsep)*\real{0.12}}@{}}
\caption*{\textbf{Table D.1.} Submitted model files}\\
\toprule
\textbf{Path} & \textbf{Content} & \textbf{Size} \\
\midrule
\endfirsthead
\toprule
\textbf{Path} & \textbf{Content} & \textbf{Size} \\
\midrule
\endhead
\path{task2/models/lr/} & Spark ML pipeline: logistic regression, with its input transformations & 0.1 MB \\
\path{task2/models/rf/} & Spark ML pipeline: random forest, 50 trees & 49 MB \\
\path{task2/models/mapper.json} & threshold mapper: how the rain of each class is spread, used to turn class probabilities into P(Y \(>\) \(\theta\)) & 8 KB \\
\bottomrule
\end{longtable}
\endgroup

To use them, with the feature tables of \path{task2/data/} in place:

\begin{verbatim}
python task2/code/predict.py S90 "2024-12-29 16:40" 5
python task2/code/predict.py S90 "2024-12-29 16:40" 5 --model LR
\end{verbatim}

